\section{Abel}

After clearing the denominators to ensure all algebraic expressions are expressed as polynomials in the radicals, Niels Henrik Abel's 1826 proof proceeds through a structured combination of field extensions, 
permutation group properties, and a final proof by contradiction.
Abel did not use modern Galois group vocabulary, but his strategy directly mapped out what we now call the Galois theory of fields.\\

{\bf 1. The Structure of an Algebraic Solution}
Abel begins by formalizing what a general ``solution by radicals'' must look like. 
If a general quintic equation is solvable, its roots must be expressible by a finite chain of algebraic operations. 
He structures the general form of a root as a polynomial in a radical:\(x=q_{0}+q_{1}R^{1/m}+q_{2}R^{2/m}+\dots +q_{m-1}R^{(m-1)/m}\)
Here, \(m\) is a prime number, and the quantities \(q_{i}\) and \(R\) are either rational functions of the original coefficients, or they belong to a lower "order" field extension previously built up by adjoining other radicals.\\

{\bf 2. Formulating Abel's Theorem on Radical Functions}
The core pivot of the proof relies on a fundamental theorem Abel establishes about the algebraic dependencies of these radicals. 
He proves that if a root of a general quintic can be written in the form above, then the radical itself,
 \(R^{1/m}\), must be expressible as a rational function of the roots of the quintic and the coefficients.
 By applying permutations to the roots of the general quintic, Abel studies how the value of \(R^{1/m}\) behaves. 
 Because the coefficients of the general quintic are independent, independent variables, permuting the roots causes \(R^{1/m}\) to cycle through its \(m\) distinct algebraic values (the \(m\)-th roots of unity multiplied by the radical).\\

 {\bf 3. Restricting the Prime Index \(m\)}
 Abel analyzes what values the prime degree \(m\) of the outermost radical can take.A function of 5 independent variables can only take a specific number of distinct values when its variables are permuted.
 By invoking Cauchy's earlier work on the number of values a symmetric or partially symmetric function can take, Abel shows that the number of values a rational function of 5 variables can take under permutation 
 must be a divisor of \(5! = 120\).
 He narrows down the possible values of the prime index \(m\). 
 The only primes that can mathematically divide these algebraic combinations for 5 variables are \(m = 2\) and \(m = 5\).\\

 {\bf 4. Bounding the Value Multiplicity (The Cauchy Constraint)}Abel 
 then looks at a critical theorem proved by Augustin-Louis Cauchy: a non-symmetric function of 5 variables cannot take fewer than 5 distinct values unless it takes exactly 2 values (which corresponds to the alternating group \(A_{5}\), generated by the square root of the discriminant).
 This creates a rigid algebraic trap:If \(m = 5\), the radical expression must take exactly 5 distinct values under root permutations.Abel shows that the only way a function of 5 variables can take exactly 5 values under permutation is if it behaves like a specific linear combination of the roots (similar to a Lagrange resolvent).\\

 {\bf 5. The Final Contradiction}
 Abel explicitly sets up the equations for these 5-valued functions. He assumes a root looks like:\(x=q_{0}+q_{1}R^{1/5}+q_{2}R^{2/5}+q_{3}R^{3/5}+q_{4}R^{4/5}\)
 By examining the algebraic behavior of the individual terms under 3-cycle and 5-cycle permutations, he derives a system of equations for the components \(q_{i}\).
 He proves that for the general quintic, it is impossible for these equations to hold simultaneously. 
 The permutations of the 5 independent roots require the individual components to vanish or collide in a way that contradicts the assumption that \(x\) is a root of a general equation with independent coefficients.
 Because the assumption that a general root can be expressed as a chain of radicals leads directly to an algebraic impossibility, Abel concludes that the general quintic equation cannot be solved by radicals. \\

 {\bf Summary of Abel's Path After Clearing Denominators}
 \begin{itemize}
 \item[]Defines Root Anatomy: Expresses the hypothetical root as a polynomial of a prime-degree radical \(R^{1/m}\).
 \item[]Inverts the Radical: Proves that \(R^{1/m}\) must be a rational function of the quintic's roots.
 \item[]Applies Permutations: Analyzes how \(R^{1/m}\) changes values when the roots are permuted.
 \item[]Invokes Cauchy's Bound: Restricts the possible values of \(m\) to 2 or 5 based on the combinatorics of 5 variables.
 \item[]Forces Contradiction: Demonstrates that the algebraic conditions required to form a 5-valued function cannot be met by the coefficients of a general quintic.
 \end{itemize}
 
 To understand how Augustin-Louis Cauchy's work on permutations constrained Niels Henrik Abel, 
 we have to look at the combinatorics of switching variables.\\
 When Abel proved that the outermost radical of a hypothetical solution, \(v = R^{1/m}\), must be a rational function of the five roots \((x_1, x_2, x_3, x_4, x_5)\), 
 he created a bridge between algebra and combinatorics. 
 Because the coefficients of a general quintic are completely independent, permuting the roots creates a set of distinct values for that function.\\
 Cauchy's theorem placed a strict mathematical ceiling on exactly how many values such a function could take, narrowing Abel's options down to a dead end.\\

{\bf  1. The Core of Cauchy's Theorem (1815)}
In 1815, Cauchy published a groundbreaking theorem concerning symmetric and non-symmetric functions. 
He proved a rigid restriction for functions of \(n\) independent variables
\begin{theorem}[Cauchy's Number of Values Theorem:]
If a non-symmetric function of \(n\) independent variables changes its value when the variables are permuted, the number of distinct values it can take (its index) must be at least the largest prime factor of \(n\). Furthermore, for \(n = 5\), a function cannot have 3 or 4 distinct values. 
It can only have 1 (fully symmetric), 2, 5, or more (up to \(5! = 120\)) distinct values.
\end{theorem}

{\bf 2. How This Restricted Abel's Prime Radical (\(m\))}
Abel assumed the general root of the quintic could be written as a polynomial in a radical of prime degree 
\(m\):\(x=q_{0}+q_{1}v+q_{2}v^{2}+\dots +q_{m-1}v^{m-1}\quad \text{where\ }v=R^{1/m}\)
By definition, if you multiply \(v\) by the \(m\)-th root of unity (\(\omega \)), its \(m\)-th power \(R\) stays exactly the same. 
Therefore, the function \(v(x_1, x_2, x_3, x_4, x_5)\) can take exactly \(m\) distinct values as its variables are permuted.
Abel applied Cauchy's theorem to this \(m\):Because \(v\) takes exactly \(m\) values, 
Cauchy's theorem dictates that \(m\) must be a possible number of values for a 5-variable function.
According to Cauchy, for 5 variables, the number of values can only be 1, 2, 5, or \(\ge 6\).
Since \(m\) must be a prime number by algebraic design, the only prime numbers allowed in that set are \(m = 2\) and \(m = 5\).
Cauchy's theorem instantly destroyed the possibility of using cubic roots (\(m=3\)), seventh roots (\(m=7\)), or any higher prime radicals at the final stage of a quintic solution.\\

{\bf 3. Eliminating the \(m = 2\) Option}
An \(m=2\) radical means a square root. 
If \(m=2\), the function \(v\) takes exactly 2 values under permutation.
Cauchy's work had already thoroughly mapped out 2-valued functions of 5 variables. 
The only way a function of 5 variables can take exactly 2 values is if it is a symmetric function multiplied by the alternating sign function 
(the square root of the discriminant, \(\sqrt{\Delta }\)).
Abel knew that adjoining a square root (\(\sqrt{\Delta }\)) merely splits the symmetric group \(S_{5}\) into the alternating group \(A_{5}\). 
It simplifies the problem slightly but does not solve the equation. 
A square root alone cannot express five distinct, independent roots. 
Therefore, \(m=2\) could not be the final step.\\

{\bf 4. Trapping the \(m = 5\) Option}
This left Abel with exactly one remaining choice: \(m = 5\). 
The outermost radical had to be a fifth root, meaning the function \(v\) must take exactly 5 distinct values when the roots are permuted.
Abel turned back to Cauchy's structural analysis of 5-valued functions. 
Cauchy had proven that if a function of 5 variables has exactly 5 values, the permutations that leave the function unchanged must form a specific subgroup of order 24 
(isomorphic to the symmetric group \(S_{4}\)).
To satisfy this rigid combinatorial structure, Abel showed that the function \(v\) must be a linear combination of the roots structured precisely like a 
Lagrange resolvent:\(v=x_{1}+\alpha x_{2}+\alpha ^{2}x_{3}+\alpha ^{3}x_{4}+\alpha ^{4}x_{5}\)(where \(\alpha \) is a fifth root of unity).\\

{\bf 5. The Closing of the Trap}
By using Cauchy's theorem to force \(v\) into this exact, unyielding algebraic shape, Abel was able to execute his final blow.
He applied a simple 3-cycle permutation (e.g., swapping \(x_1 \to x_2 \to x_3 \to x_1\)) to the roots. 
Because a 3-cycle permutation belongs to the alternating group \(A_{5}\) but not to the subgroup that keeps a 5-valued function stable, it forced a contradiction. 
Abel proved that under these required permutations, the algebraic coefficients would have to simultaneously equal completely different values, 
meaning the assumed fifth-root function could not actually exist for independent variables.
Without Cauchy's theorem bounding the permissible number of values to 2 and 5, Abel would have had to fight an infinite number of potential radical degrees (\(m=3, 5, 7, 11\dots\)). 
Cauchy provided the boundaries that allowed Abel to corner the quintic and prove it impossible.\\

In his 1826 paper, Niels Henrik Abel used a beautifully simple property of 3-cycle permutations to snap his algebraic trap shut.
He didn't use modern group theory notation (like \(A_{5}\) or \(S_{5}\)), but he expressed the exact same logical mechanism. 
Here is the exact step-by-step breakdown of how a single 3-cycle permutation created his final, fatal contradiction.\\

{\bf 1. The Setup:} 
The Five-Valued RadicalThrough Cauchy's theorem, Abel established that the final, outermost radical needed to solve the quintic must be a fifth root, which we can call \(v\):\(v=R^{1/5}\)
Because \(v\) is a fifth root, permuting the roots of the quintic can only change \(v\) by multiplying it by some fifth root of unity, \(\omega \) (where \(\omega^5 = 1\)). 
Therefore, under the full group of permutations, \(v\) can take exactly 5 distinct values:\(\{v,\omega v,\omega ^{2}v,\omega ^{3}v,\omega ^{4}v\}\) \\

{\bf 2. Applying the 3-Cycle Permutation}
Abel then chose an arbitrary 3-cycle permutation of the roots, which we can denote as \(p\). 
A 3-cycle simply takes three roots and rotates them in a loop (for example, swapping \(x_1 \to x_2 \to x_3 \to x_1\)), 
leaving the other two roots completely untouched.
Let's apply this 3-cycle \(p\) to our radical function \(v\). 
As established, applying any permutation to the roots can only multiply \(v\) by some fifth root of unity, say \(\omega ^{k}\):\(p(v)=\omega ^{k}\cdot v\) \\

{\bf 3. The Power of Order Three}
A 3-cycle has a strict geometric property: if you execute it exactly three times, every element returns to its original starting position (\(p^3 = e\), the identity permutation).
Abel applied the exact same 3-cycle permutation three times in a row to both sides of the equation:
\begin{itemize}
\item[]First time: \(p(v) = \omega^k \cdot v\)
\item[]Second time: \(p(p(v)) = \omega^k \cdot \omega^k \cdot v = \omega^{2k} \cdot v\)
\item[]Third time: \(p(p(p(v))) = \omega^k \cdot \omega^k \cdot \omega^k \cdot v = \omega^{3k} \cdot v\)
\end{itemize}
Because applying a 3-cycle three times restores the roots to their original order, the left side must evaluate back to the original, un-permuted value of \(v\):\(v=\omega ^{3k}\cdot v\)
For this equation to hold true, the multiplier must equal 1:\(\omega ^{3k}=1\) \\

{\bf 4. The Arithmetic Lock}
We now have two strict rules dictating the behavior of this multiplier:From the nature of the fifth root, \(\omega \) is a fifth root of unity (\(\omega^5 = 1\)).
From the nature of the 3-cycle, \(\omega ^{3k}\) must equal 1.
Because 3 and 5 are prime numbers that share no common factors (coprime), the only power that satisfies both conditions simultaneously is \(k = 0\), meaning \(\omega^0 = 1\).
Therefore:\(p(v)=1\cdot v=v\)
This reveals an astonishing fact: every single 3-cycle permutation leaves the radical \(v\) completely unchanged.\\ 

{\bf 5. The Final Contradiction}
This is where Abel applies the coup de gr\"ace using Cauchy's framework:\\
\begin{itemize}
\item[]Fact A: In combinatorics, the 3-cycle permutations are the foundational building blocks that generate the entire {\bf Alternating Group (\(A_{5}\))}. 
Because \(v\) is unaffected by every 3-cycle, it is completely invariant under the entire alternating group.
\item[]Fact B: According to Cauchy's structural theorems, any function of 5 variables that is completely invariant under the alternating group can take {\bf at most 2 distinct values} across all possible permutations (acting like the square root of a discriminant).
\end{itemize}
Now, look at the two numbers we have forced upon the exact same function \(v\): From its definition as a fifth root, \(v\) must take exactly 5 values.
From the 3-cycle arithmetic, \(v\) can take at most 2 values.
A function cannot simultaneously possess exactly 5 values and at most 2 values. 
The assumption that a general quintic can be solved by a chain of radicals containing a fifth root collapses into an absolute algebraic impossibility.\\
